\section{Claim boundary and operational consequences}
\label{sec:claim-boundary}

\begin{theorem}[Audited conclusion]
\label{thm:audited-conclusion}
For every finitely supported complex coefficient,
\[
 \Flat=\canc^2\Seff.
\]
For every nonzero coefficient,
\[
 (2-\Seff)_+\le\Flat\le\Seff\le S.
\]
For the literal TPC-32 determinant coefficient at one fixed nonzero
\(h_0\), natural post-bin mass survival implies
\[
 \SeffX=X^{o(1)}Q_X.
\]
At
\[
 \beta=\frac{267}{400},
\]
the TPC-32 condition
\[
 \Flat_X\le X^{1/400+o(1)}
\]
is then equivalent to
\[
 \frac{|Z_X|}{M_X}
 \le\frac{X^{o(1)}}{J_X},
\]
and to the absolute target
\[
 |Z_X|\le X^{o(1)}Q_X^2.
\]
\end{theorem}

\begin{proof}
Combine
\cref{thm:exact-factorization,thm:sharp-strip,thm:natural-effective-support,thm:critical-one-over-J,cor:relative-absolute}.
\end{proof}

\paragraph{What is proved.}
The factorization, sharp strip, equality cases, and finite
countermodels are unconditional \(\Lzero\) results.  Their
identification on the already aggregated fixed-\(h_0\) coefficient,
together with the conditional natural-mass and exponent transfers,
is an \(\Lone\) result.

\paragraph{What is not proved.}
This paper proves no lower bound for \(M_X\) or \(D_X\), no upper
bound for \(\canc_X\) or \(Z_X\), and no new fixed-shift
M{\"o}bius-correlation estimate.  It proves no parity breakthrough, no
Hardy--Littlewood asymptotic, no prime-pair lower bound, and no
twin-prime result.  Nothing is specialized from a general nonzero
\(h_0\) to \(h_0=2\).

\paragraph{Edge-case rule.}
If \(A_X=0\), then \(M_X=D_X=Z_X=0\) and the convention
\(\canc_X=\SeffX=\Flat_X=0\) preserves the exact factorization and
the nonnegative upper bounds.  The lower strip
\((2-\SeffX)_+\le\Flat_X\) is asserted only for \(A_X\ne0\).
This trivial flatness does not certify natural mass or energy
survival.  No expression containing \(M_X^{-1}\) or \(D_X^{-1}\)
may be used on that branch.

\paragraph{Normalization rule.}
Only the post-bin mass
\[
 M_X=\sum_n|A_X(n)|
\]
enters the factorization.  A raw absolute mass or TPC-32 majorant
cannot be substituted for it.  Likewise, actual support
\(S_X\asymp Q_X\) would not by itself imply
\(\SeffX\asymp Q_X\).

\paragraph{Exponent rule.}
The shorthand \(1/J_X\) is valid at the critical endpoint only in
the natural mass--energy corner.  With nonzero losses, the correct
relative requirement is
\[
 2\eta_{\canc}
 \ge
 \beta-\chi+\lambda-2\mu.
\]
At \(\chi=1/400\) this is equivalent to the absolute ledger
\(2\eta_Z\ge\lambda\).  Neither the geometric \(1/400\) baseline nor
any mass exponent may be spent twice.

\paragraph{Independent physical-loss rule.}
The compatibility \(2\eta_Z\ge\lambda\) is not the full-route loss
budget.  Separately, every fixed-power loss actually entering the
physical interface and downstream comparison, but not already
represented in \(\lambda-2\eta_Z\), must satisfy
\[
 \Lambda_{\mathrm{phys}}<\frac1{400}.
\]
Equality or excess stops the route, and no loss is charged in both
ledgers.

\paragraph{Next gate.}
After natural post-bin mass is established, the remaining arithmetic
task is exactly a relative phase-cancellation theorem at scale
\[
 \frac{|Z_X|}{M_X}\ll\frac{X^{o(1)}}{J_X}.
\]
Without natural mass, the full profile
\((\mu,\lambda,\eta_{\canc})\) must be measured first.  A numerical
observation, averaged-shift result, or generic-frequency estimate may
motivate that theorem but is not an \(\Ltwo\) substitute for the
literal fixed-\(h_0\) bound.
